\documentclass[11pt]{article}
\usepackage{fontspec}
%\usepackage[review]{acl}
\usepackage[round]{natbib}
\usepackage[margin=0.5in]{geometry}
\usepackage{array}
\usepackage{booktabs}
\usepackage[normalem]{ulem}
\usepackage[colorlinks=true, urlcolor=blue, linkcolor=blue]{hyperref}
\usepackage{float}
\setmainfont{Arial}
\title{Detecting Machine Generated Texts in isiZulu \\[1ex] \large Evaluating pre-trained models and diverse datasets}
\author{Ruan Esterhuizen (u23532387), Ruben Hannes Gadd (u23633353), Christiaan Strydom (u22498878) \\[1ex] \large Group 19}
\date{}
\begin{document}
    \maketitle
    \section{Introduction}
    With the dawn of the 20th decade of the 21st millennium, the arrival of AI models such as ChatGPT has 
    caused rippling effects across the globe both positive and negative. Such negative effects include 
    the abuse of language models or "chatbots" to generate text content that is difficult to differentiate 
    from text written by man. With further advancements into language models over the years, it has become 
    more challenging for humans to tell whether a piece of text is machine-generated or not. This research 
    report aims to provide a solution to the problem via the utilization of NLP techniques to classify if 
    a given text is MGT (Machine generated Text) or HGT (Human Generated Text). The research report details 
    the investigation into creating a Zulu-based MGT detection tool. Zulu/isiZulu was chosen as our primary 
    language since it is the most spoken African language in South Africa and shares linguistic roots with 
    other African languages that fall under the Nguni linguistic tree which include isiXhosa, siSwati and 
    isiNdebele. Additionally, since isiZulu is not as globally spoken outside of Africa, there is limited 
    number of models that prioritize generating text in isiZulu. Note that the use of abbreviations 
    MGT (Machine Generated Text) and HGT (Human Generated Text) will be used throughout the research report.

    \section{Literature Review}
    
    \section{Methodology}
    %Add small 'intro' if possible
    \subsection{Gathering Datasets}
    The datasets selected to be used for our HGT dataset are datasets 
    isiZulu News, AfriHate* and Vuk'uzenzele which are a collection 
    of news articles, social media posts and magazine articles in isiZulu. 
    Python scripts were created to retrieve the datasets and store them as 
    csv files. 
    
    \begin{table}[H]
    \centering
    \begin{tabular}{l l l l l}
    \toprule
    \textbf{Name} & \textbf{Domain} & \textbf{Size (items)} & \textbf{Description} & \textbf{Repository} \\
    \midrule
    Vuk'uzenzele 
        & Magazine articles 
        & 185 
        & Articles from SA government magazine 
        & \href{https://github.com/dsfsi/vukuzenzele-nlp}{\uline{link}} \\
    IsiZulu News 
        & News articles 
        & 752 
        & Scraped publicly available news articles
        & \href{https://github.com/dsfsi/za-isizulu-siswati-news-2022}{\uline{link}} \\
    AfriHate* 
        & Social media 
        & $\sim$4000 
        & Scraped public Twitter posts 
        & \href{https://github.com/AfriHate/AfriHate}{\uline{link}} \\
    \end{tabular}
    \caption{Datasets selected for HGT}
    \end{table}
    Bear in mind that the AfriHate* dataset contains benign yet sensitive content 
    that may be considered offensive to certain groups, yet was still selected due 
    to being a huge source of text in isiZulu.
    \subsection{Pre-processing Datasets}
    \subsubsection{isiZulu News}
    \subsection{Generating MGT}
    \section{Experiments and Results}

    \section{Reflections}

    \section{Conclusion}
    %\bibliographystyle{plainnat}
    %\bibliography{reportReferences}
\end{document}